\documentclass[11pt]{article}
\usepackage[letterpaper,margin=1.1in]{geometry}
\usepackage[T1]{fontenc}
\usepackage{lmodern}
\usepackage{microtype}
\usepackage{amsmath}
\usepackage{booktabs}
\usepackage{array}
\usepackage[round,authoryear]{natbib}
\usepackage{hyperref}
\hypersetup{colorlinks=true,linkcolor=black,citecolor=black,urlcolor=black,
  pdftitle={AutoCAD-Compatible Unit Scaling in the GoSurvey Traverse Editor and Clipboard},
  pdfauthor={Chet Jones}}
\usepackage{enumitem}
\setlength{\parindent}{0pt}
\setlength{\parskip}{0.55\baselineskip}
\newcolumntype{L}[1]{>{\raggedright\arraybackslash}p{#1}}

\title{\textbf{Is a Drawing Unit Just a Label?\\AutoCAD-Compatible Unit Scaling in\\GoSurvey's Traverse Editor and Clipboard}}
\author{Chet Jones \\ \small GoSurvey Project}
\date{October 2026}

\begin{document}
\maketitle
\vspace{-1.5em}
\begin{center}
\small REQ-357 \quad\textbar\quad REQ-383 \quad\textbar\quad GitHub Issue \#831 \quad\textbar\quad Pull Request \#837
\end{center}
\vspace{1em}

\begin{abstract}
\noindent A Reddit report that the Traverse Closure Analysis panel labelled
distances \texttt{(ft)} on a meters project led to a harder question: if a
drawing's linear unit is just a label on a number, does changing it --- or
moving a number between two drawings with different labels --- ever change
anything real? This note answers both halves precisely. First, it confirms
by inspection of GoSurvey's own code that changing a drawing's unit setting
never rescales existing geometry, which is not a bug but the same
convention AutoCAD's own \texttt{INSUNITS} system variable uses
\citep{autodeskInsunits}. Second, it finds and closes a real gap against
that same convention: AutoCAD scales content pasted or inserted between two
drawings of different units by the ratio of those units
\citep{autodeskInsunits}, but GoSurvey's clipboard paste skipped that check
entirely whenever neither drawing belonged to one of its own ``Projects''.
This note derives the governing unit-ratio formula, works a numeric example
by hand, and documents the warn-then-convert implementation and its test
suite --- deliberately never silent, unlike AutoCAD's own default.
\end{abstract}

\section{The question: is the unit label load-bearing?}
\label{sec:question}

GoSurvey stores every drawn length as a plain number with no unit attached
--- a line ``100 units long'' is stored as the double \texttt{100.0}, full
stop. The drawing's \texttt{INSUNITS} code (Unitless, Inches, Feet,
Millimeters or Meters; REQ-357) is a separate setting that says how that
number should be read: 100 feet, or 100 meters, or something else entirely.
Changing the setting after the fact does \emph{not} walk back through the
drawing and rescale every stored number to the new unit --- the number
\texttt{100.0} stays \texttt{100.0}, now read under the new label.

This is confirmed directly from the source rather than assumed: every
write site for the drawing's unit setting
(\texttt{src/ui/CadUiSettings.cpp}, \texttt{src/ui/CadUi\_DrawingSettings.cpp})
assigns only the integer code itself (\texttt{cmd.drawingInsUnits = ...}); none
of them touches a coordinate, a distance or any entity array. No code path
in the program rescales geometry in response to a unit-setting change.

\emph{This is deliberate, not an oversight, and it is the same rule
AutoCAD follows.} The Autodesk reference for \texttt{INSUNITS} describes it
as the value AutoCAD uses to \emph{auto-scale} content \textbf{being
inserted or attached} --- blocks, images, and xrefs --- into a drawing; it
does not describe \texttt{INSUNITS} as retroactively rescaling anything
already drawn \citep{autodeskInsunits}. Flipping the unit dropdown on an
existing drawing, in both programs, only changes how future insertions and
pastes are scaled and how existing numbers are labelled on screen --- never
what those numbers already are.

\section{The real gap: paste between two differently-labelled drawings}
\label{sec:gap}

If the label is informational for an existing drawing, the question moves
to the boundary where two drawings meet: what happens when you copy an
object out of a meters drawing and paste it into a feet drawing?

\subsection*{What AutoCAD does}

AutoCAD computes the ratio of the source and destination drawings'
\texttt{INSUNITS} values and scales the inserted or pasted content by that
ratio, so that the object keeps its real-world size rather than its raw
number \citep{autodeskInsunits}. A millimeter-drawn object pasted into an
inch-based drawing is scaled by $25.4$, the millimeters-per-inch ratio, as a
direct worked instance of the same rule \citep{autodeskInsunits}. This
scaling is automatic and, by default, silent: nothing stops to ask first.
AutoCAD does expose a toggle for it (\texttt{DWGUNITS}'s ``Scale objects
from other drawings upon insert''), and user reports of surprising scale
jumps are common enough that the standard community advice is simply to
turn it off and keep every file in one unit \citep{autodeskInsunits}.

\subsection*{What GoSurvey did, before this change}

GoSurvey already replicated half of this correctly: inserting a
\textsc{block} compares the block's own recorded unit against the current
drawing's unit and computes a scale factor automatically
(\texttt{CadBlockInsertUnitsScale}, gated by the existing REQ-357
``Scale inserted objects'' toggle --- GoSurvey's own equivalent of
AutoCAD's \texttt{DWGUNITS} option).

Clipboard \textsc{paste}, however, did not. \texttt{CheckClipboardPaste}
(\texttt{src/commands/ProjectWarnings.cpp}) returned immediately, with no
unit check performed at all, whenever neither the source nor the
destination drawing belonged to a GoSurvey \emph{Project} (its own
survey-point-sharing feature, unrelated to AutoCAD). REQ-383 clause 7 had
documented this explicitly: ``a paste between two standalone drawings is
unchanged.'' Pasting a $100.0000$-unit object from a meters drawing into a
feet drawing landed it as $100.0000$ feet --- a real $228\%$ oversize,
silently, because no unit arithmetic ran at all.

\section{The governing formula}
\label{sec:formula}

GoSurvey already carries a meters-per-unit table for its five
\texttt{INSUNITS} codes (\texttt{geo::MetersPerInsUnit}):
\begin{center}
\begin{tabular}{@{}lcl@{}}
\toprule
Code & Unit & Meters per unit \\
\midrule
1 & Inches & $0.0254$ \\
2 & Feet (international) & $0.3048$ \\
2 & Feet (U.S.\ survey) & $1200/3937 = 0.3048006096\ldots$ \\
4 & Millimeters & $0.001$ \\
6 & Meters & $1.0$ \\
\bottomrule
\end{tabular}
\end{center}
When only the unit differs (no map-projection zone is involved), the
transform GoSurvey's existing conversion engine computes is a pure uniform
scale about the world origin \citep{goSurveyRepo}:
\begin{equation}
k_{\mathrm{unit}} \;=\; \frac{m_{\mathrm{from}}}{m_{\mathrm{to}}},
\label{eq:kunit}
\end{equation}
where $m_{\mathrm{from}}$ and $m_{\mathrm{to}}$ are the source and
destination drawings' meters-per-unit values (\texttt{geo::PlanConversion},
the \texttt{!plan.zoneDiffers} branch). Every pasted coordinate, length and
height is multiplied by $k_{\mathrm{unit}}$:
\begin{equation}
x_{\mathrm{dest}} = k_{\mathrm{unit}} \cdot x_{\mathrm{src}}, \qquad
y_{\mathrm{dest}} = k_{\mathrm{unit}} \cdot y_{\mathrm{src}}, \qquad
z_{\mathrm{dest}} = k_{\mathrm{unit}} \cdot z_{\mathrm{src}}.
\label{eq:scaleapply}
\end{equation}
Equation~\ref{eq:kunit} is exactly the same ratio AutoCAD computes for a
cross-unit insert or paste \citep{autodeskInsunits}; GoSurvey reuses its
existing conversion engine rather than adding a second formula.

\subsection*{Worked example: meters into feet}

A $100.0000$-unit line, drawn in a meters drawing ($m_{\mathrm{from}} =
1.0$), copied and pasted into a feet drawing ($m_{\mathrm{to}} = 0.3048$,
international foot):
\begin{center}
\begin{tabular}{@{}llr@{}}
\toprule
Quantity & Formula & Value \\
\midrule
Source length & drawn & $100.0000$ (m) \\
$k_{\mathrm{unit}}$ & $m_{\mathrm{from}} / m_{\mathrm{to}} = 1.0 / 0.3048$ & $3.2808398950$ \\
Pasted length, unconverted (before this change) & no scaling applied & $100.0000$ ft \\
Pasted length, converted (after this change) & $100.0000 \times k_{\mathrm{unit}}$ & $328.0839895$ ft \\
\bottomrule
\end{tabular}
\end{center}
A $100$ m line is $328.0839895$ ft by definition ($1$ m $= 1/0.3048$ ft
exactly); the converted figure matches to the number of decimals shown,
confirming Equation~\ref{eq:scaleapply} recovers the physically correct
length rather than an approximation.

\section{What this change does, and deliberately does not do}
\label{sec:implementation}

\texttt{CheckClipboardPaste} (\texttt{src/commands/ProjectWarnings.cpp}) now
computes the unit mismatch check \emph{before} the standalone-drawing early
return, rather than inside the branch that early return used to skip
entirely. When both drawings are standalone and their units differ, the
function returns \texttt{PasteCheck::Verdict::Warn} with an explanatory
message and \texttt{canConvert} set by the same
\texttt{PlanClipboardConversion} / \texttt{ApplyClipboardConversion} pair
REQ-383 already uses for a Project-linked units mismatch --- no second
conversion path was written. The existing paste-prompt UI
(\texttt{st.pastePrompt}) already handled a non-blocking warning with a
convert option generically, so no new UI code was needed either.

\textbf{Deliberately unchanged:} a coordinate-system (map-projection zone)
mismatch between two standalone drawings is still not checked. AutoCAD has
no equivalent concept to check against for a plain drawing pairing, and
extending that check would require deciding what a ``zone'' even means for
two arbitrary standalone drawings --- a materially different, unasked-for
question from the one this change answers. Equation~\ref{eq:kunit}'s
scope is unit ratio only.

\textbf{Deliberately not silent:} AutoCAD applies this scaling without
asking, which its own user community reports as a frequent source of
surprise \citep{autodeskInsunits}. GoSurvey's version always surfaces the
warning and lets the user choose \emph{Convert} or \emph{Cancel} ---
replicating the arithmetic, not the silence.

\section{Automated test summary}

\begin{center}
\small
\begin{tabular}{@{}L{0.3\textwidth}L{0.34\textwidth}L{0.28\textwidth}@{}}
\toprule
File & Test case & What it proves \\
\midrule
\texttt{ProjectWarningsTests.cpp} &
  two standalone drawings and the same project are never checked &
  a ZONE-only mismatch between standalone drawings still returns \texttt{Ok} (unchanged) \\[0.3em]
\texttt{ProjectWarningsTests.cpp} &
  two standalone drawings with a unit mismatch still warn and offer convert (issue \#831) &
  a UNIT-only mismatch between standalone drawings now returns \texttt{Warn} with \texttt{canConvert}, and \texttt{StartPasteCommand} opens a non-blocking prompt \\[0.3em]
\texttt{ProjectWarningsTests.cpp} &
  a units mismatch warns / can be converted: lengths, heights, radii and the origin all follow &
  pre-existing Project-linked path is unchanged by this edit \\
\bottomrule
\end{tabular}
\end{center}

All cases pass under the \texttt{GoSurveySnapTests} executable (Catch2,
MSVC release build); the four unrelated pre-existing \texttt{[req383]}
failures in this build (point-visibility renumbering, tracked separately)
are identical before and after this change, confirmed by running the same
filtered suite against both revisions. The source is PR~\#837, branch
\texttt{fix/issue831-traverse-closure-units}, amending REQ-383 clause 7 in
\texttt{spec/requirements.md} (decision \texttt{D-2026-10-11-c} in
\texttt{spec/project.md}).

\begin{thebibliography}{9}

\bibitem[Autodesk, 2026]{autodeskInsunits}
Autodesk, Inc.\ (retrieved 2026). \emph{INSUNITS (System Variable)},
AutoCAD / AutoCAD~LT Help. Autodesk Knowledge Network.
\url{https://help.autodesk.com/cloudhelp/2022/ENU/AutoCAD-LT-MAC/files/GUID-A58A87BB-482B-4042-A00A-EEF55A2B4FD8.htm}

\bibitem[GoSurvey, 2026]{goSurveyRepo}
GoSurvey Project (2026). \emph{Source repository}, functions
\texttt{geo::MetersPerInsUnit}, \texttt{geo::PlanConversion}
(\texttt{src/geo/DrawingConversion.cpp}), \texttt{CheckClipboardPaste}
(\texttt{src/commands/ProjectWarnings.cpp}), and
\texttt{CadBlockInsertUnitsScale} (\texttt{src/commands/CadBlocks.cpp}).
\url{https://github.com/chetjones003/GoSurvey}

\end{thebibliography}

\end{document}
